\documentclass[11pt,a4paper]{article}
\usepackage[margin=23mm,headheight=15pt]{geometry}
\usepackage{fontspec}
\newfontfamily\codemathfont{latinmodern-math.otf}[Path=fonts/]
\newcommand{\CodeSymbol}[1]{\ifnum"#1="29F8\ensuremath{\big/}\else{\codemathfont\symbol{"#1}}\fi}
\usepackage{amsmath}
\usepackage{unicode-math}
\usepackage{xcolor,booktabs,longtable,array,enumitem,microtype}
\PassOptionsToPackage{obeyspaces}{url}
\usepackage{listings,fvextra,xurl,hyperref,bookmark,fancyhdr,titlesec,needspace,caption}
\definecolor{teal}{HTML}{164C51}
\definecolor{muted}{HTML}{516466}
\definecolor{pale}{HTML}{F2F6F6}
\hypersetup{colorlinks=true,linkcolor=teal,urlcolor=teal,pdftitle={Paper II: Lean Guide},pdfauthor={Komorebi P.},bookmarksnumbered=true}
\urlstyle{tt}
\newcommand{\code}[1]{\nolinkurl{#1}}
\newcommand{\source}[1]{\par\smallskip{\footnotesize\color{muted}Source: \code{#1}.}\par\smallskip}
\newcommand{\excerpt}[2]{\Needspace{5\baselineskip}\source{KomorebiVerification/#1.lean}\VerbatimInput[fontsize=\small,frame=leftline,rulecolor=\color{teal},framesep=6pt,breaklines=true,breakanywhere=true,commandchars=\\\{\},xleftmargin=8pt]{excerpts/#2.render.txt}}
\titleformat{\section}{\Large\sffamily\bfseries\color{teal}}{\thesection}{.65em}{}
\titleformat{\subsection}{\large\sffamily\bfseries\color{teal}}{\thesubsection}{.65em}{}
\titleformat{\subsubsection}{\normalsize\sffamily\bfseries\color{teal}}{\thesubsubsection}{.65em}{}
\titlespacing*{\section}{0pt}{2.2ex plus .6ex}{1ex}
\titlespacing*{\subsection}{0pt}{2ex plus .4ex}{.7ex}
\setlength{\parindent}{0pt}
\setlength{\parskip}{.55em}
\setlength{\emergencystretch}{2em}
\setlist{nosep,topsep=.4em,leftmargin=1.5em}
\pagestyle{fancy}\fancyhf{}
\fancyhead[L]{\small\sffamily\color{muted}KOMOREBI \quad PAPER II}
\fancyhead[R]{\small\sffamily\color{muted}LEAN GUIDE}
\fancyfoot[C]{\small\thepage}
\renewcommand{\headrulewidth}{.3pt}
\lstset{basicstyle=\ttfamily\small,backgroundcolor=\color{pale},frame=single,rulecolor=\color{pale},framerule=0pt,xleftmargin=5pt,xrightmargin=5pt,aboveskip=7pt,belowskip=7pt,breaklines=true,breakatwhitespace=false,columns=fullflexible,keepspaces=true,showstringspaces=false,tabsize=2,escapeinside={(*@}{@*)}}
\makeatletter
\renewcommand*\l@subsection{\@dottedtocline{2}{1.5em}{3.5em}}
\makeatother

\usepackage{komorebi-monograph}
\renewcommand{\KomorebiKind}{Lean and evidence guide}
\renewcommand{\KomorebiCoverBlurb}{Current formal scope and evidence, followed by the inherited declaration-level walkthrough and exact source excerpts.}
\setmonofont{DejaVuSansMono}[Path=fonts/,UprightFont=*.ttf,BoldFont=*-Bold.ttf,ItalicFont=*-Oblique.ttf,BoldItalicFont=*-BoldOblique.ttf,Scale=.80]
\input{generated/current_edition.tex}
\KomorebiSetup
\begin{document}
\KomorebiCover
\KomorebiOpeningNote{This guide separates the current formal development from the inherited 25-row contract. The current overview records exact source scope and verification status. The later walkthrough preserves the canonical source architecture, representative proofs, and historical evidence with their original boundaries.
\medskip
\input{generated/current_opening.tex}}


\tableofcontents
\clearpage
\pagenumbering{arabic}
\input{generated/current_release.tex}
\clearpage
\input{chapters/07_analytic_closeout.tex}
\clearpage
\input{chapters/01_orientation.tex}
\input{chapters/02_refinement.tex}
\input{chapters/03_transfer.tex}
\input{chapters/04_arithmetic.tex}
\input{chapters/05_defects.tex}
\input{chapters/06_replay.tex}
\clearpage
\section{The inherited contract of 25 rows}\label{sec:map}
The following map preserves the canonical v1.1.1 classifications, checked scopes, and explicit boundaries. The targets use the namespace abbreviations below. These abbreviations are only typography: the companion JSON expands every name and records its source file and line. A declaration may occur in several rows. There are 171 occurrences and 160 distinct names. Definitions, structures, equivalences, and theorems all count as declaration targets.
\input{generated/namespace_prefixes.tex}
\input{generated/theorem_map.tex}
\clearpage
\appendix
\section{Navigation in the final traditional paper}\label{app:crosswalk}
The printed page numbers below refer to the final traditional companion, whose cover is unnumbered and whose preliminary pages use Roman numerals. They are not PDF viewer page indices. A shared label is a navigation aid. It does not extend the formal scope to additional claims in the surrounding section. In particular, the Smith phase display has no equation number.
\begin{longtable}{p{.07\textwidth}p{.50\textwidth}p{.32\textwidth}}
\toprule Row & Canonical label & Final locator \\ \midrule\endhead
1 & \code{def:refinement-map} & Statement 3.2, p.~3 \\ \addlinespace
2 & \code{thm:nested-refinement} & Statement 3.3, p.~3 \\ \addlinespace
3 & \code{def:coordinate-refinement} & Statement 3.4, p.~4 \\ \addlinespace
4 & \code{prop:perfect-reconstruction} & Statement 3.5, p.~4 \\ \addlinespace
5 & \code{thm:transfer-block-refinement} & Statement 3.8, p.~6 \\ \addlinespace
6 & \code{prop:residue-algebra} & Statement 4.2, p.~9 \\ \addlinespace
7 & \code{thm:root-of-unity} & Statement 4.3, p.~10 \\ \addlinespace
8 & \code{thm:polyphase-refinement} & Statement 4.6, p.~11 \\ \addlinespace
9 & \code{thm:dyadic-odd} & Statement 4.9, p.~11 \\ \addlinespace
10 & \code{thm:weighted-inverse} & Statement 5.10, p.~17 \\ \addlinespace
11 & \code{cor:finite-precision} & Statement 5.11, p.~18 \\ \addlinespace
12 & \code{thm:transfer-zeta} & Statement 6.6, p.~21 \\ \addlinespace
13 & \code{cor:zeta-recurrence} & Statement 6.7, p.~22 \\ \addlinespace
14 & \code{thm:bezout-projector} & Statement 7.5, p.~25 \\ \addlinespace
15 & \code{def:separation-weight} & Statement 7.6, p.~26 \\ \addlinespace
16 & \code{thm:precision-delay} & Statement 7.7, p.~26 \\ \addlinespace
17 & \code{prop:dyadic-sign} & Statement 9.3, p.~32 \\ \addlinespace
18 & \code{thm:smith-solve} & Statement 10.6, p.~36 \\ \addlinespace
19 & \code{thm:power-defect} & Statement 11.1, p.~37 \\ \addlinespace
20 & \code{cor:defect-norm} & Statement 11.3, p.~38 \\ \addlinespace
21 & \code{thm:resolvent-defect} & Statement 11.5, p.~38 \\ \addlinespace
22 & \code{eq:smith-inversion} & Section 12.3, Smith-phase identity, p.~43 \\ \addlinespace
23 & \code{alg:residue} & Statement B.1, p.~47 \\ \addlinespace
24 & \code{alg:weighted-inverse} & Statement B.2, p.~47 \\ \addlinespace
25 & \code{alg:path-zeta} & Statement B.3, p.~48 \\ \addlinespace
\bottomrule\end{longtable}

\section{Source imports and pinned dependencies}\label{app:imports}
This table records direct local imports. It is a graph of file dependencies rather than a logical dependency analysis of the theorems. The JSON companion also retains every direct external import. Some modules import the broad \code{Mathlib} module, while others list focused imports.
\input{generated/local_imports.tex}
The exact toolchain string is \code{leanprover/lean4:v4.34.0-rc2}. The Lake manifest pins the following source revisions. Its descriptive \code{inputRev} fields such as \code{master} and \code{main} do not replace these exact revisions.
\input{generated/dependency_pins.tex}
\section{Source excerpts and reproduction inventory}\label{app:inventory}
Every printed Lean listing comes from an exact UTF-8 excerpt in \code{excerpts/*.lean.txt}. The adjacent \code{*.render.txt} file escapes TeX control characters and substitutes a math font glyph where the code font lacks one. It preserves token order. Exact source text, not the display fallback, is the object to copy into an editor. Excerpts retain their surrounding source context only when that context is printed. They are not standalone Lean files.
\code{generated/excerpts.json} records each source path, inclusive source line range, the SHA-256 digests of each full source file, exact excerpt, and display file. \code{generated/theorem_map.json} contains all 25 rows with fully qualified targets, classifications, scopes, boundaries, and lexical source locations. \code{generated/mapped_source_index.json} provides an alphabetical target index. \code{generated/static_identity_audit.json} records all manifest comparisons and the optional baseline comparison performed for this edition.

\code{build_guide_data.py} uses only the Python standard library. It generates the map, import table, pins, excerpts, and current static audit without invoking Lean. \code{build_guide.sh} calls that generator and runs XeLaTeX four times. The bundled P052, Nimbus Sans, DejaVu Sans Mono, and Latin Modern Math fonts provide the typography for this edition. The required LaTeX packages are listed in the preamble. The independent PDF check script uses PyMuPDF and Pillow for page rendering and contact sheets.

The historical source snapshot remains preserved. Current additions live in the separate current project and are indexed in the current correspondence. Editing this guide cannot by itself repair or strengthen a Lean theorem. A mathematical change to the formal component requires a separate source change, a new compile, an updated map when appropriate, and new execution evidence.
\end{document}
